% Ch7 self-study -- "I do / You do" ladder for Zumdahl 7.1-7.13, four
% YOUR TURN questions per worked example. Wavelengths, elements and PES
% data chosen disjoint from ch07-notes and ws1-ws3.
%
% PROJECT CONVENTION ENFORCED THROUGHOUT: every trend explanation is
% Coulombic -- nuclear charge, distance, shielding. Octet / "wants to be
% stable" language earns zero on this project's rubrics, so the teaching
% text must never model it. Ladder 8 says so explicitly.
\documentclass{apchem-worksheet}

\apcunitword{Chapter}
\apcunit{7}
\apcunittitle{Atomic Structure and Periodicity}
\apcdoctitle{Self-Study \textbullet\ Chapter 7, I~do / You~do}
\apcsubtitle{Zumdahl \S7.1--7.13 \textbullet\ PDF pp.\ 331--390 \textbullet\ four YOUR TURN questions per ladder}

\begin{document}
\apctitleblock

\begin{apcnote}[How to use these notes --- read this first]
Each skill is a \textbf{ladder}: a fully worked example in the
solid-framed box, then \textbf{four} YOUR TURN questions in the dashed
box --- same skill, new numbers, no help.

\begin{enumerate}[leftmargin=*,itemsep=0.15em]
  \item \textbf{Work all four} before checking anything.
  \item Check against the gray \emph{check:} line. All four right on the
        first try earns the tick on the tracker (last page).
  \item Any misses: re-read the worked example, then redo the missed ones
        from a blank page.
\end{enumerate}

\textbf{The one sentence this chapter is built on.} Every property in
Ladders 5--8 --- orbital energy, ionization energy, atomic radius,
electronegativity --- is decided by \textbf{Coulomb's law}: how strongly
the nucleus pulls on an electron, which depends on the nuclear charge,
the distance, and how much shielding sits in between. If you can say
those three things about any electron, you can predict every trend
without memorizing a single arrow.

\textbf{Constants:} $c = \SI{3.00e8}{\meter\per\second}$,
$h = \SI{6.626e-34}{\joule\second}$,
$N_A = \SI{6.022e23}{\per\mole}$.
\end{apcnote}

% =====================================================================
\section*{\sffamily Ladder 1 \textbullet\ Wavelength, frequency, energy}
\zsec{7.1--7.2}

\[ c = \lambda\nu \qquad E = h\nu = \frac{hc}{\lambda} \]

Read the second equation as a sentence: \textbf{short wavelength means
high frequency means high energy}. Gamma rays sit at one end, radio waves
at the other, and visible light occupies a narrow band from about
\SI{400}{\nano\meter} (violet) to \SI{700}{\nano\meter} (red).

A photon's energy is per \emph{photon}. Multiply by $N_A$ to reach the
per-\emph{mole} figures chemists quote.

\begin{workedexample}[I do: one photon, three quantities]
Green light has a wavelength of \SI{525}{\nano\meter}. Find its
frequency, the energy of one photon, and the energy per mole.

\textbf{Metres first --- always convert before substituting:}
\[ 525~\text{nm} = \SI{5.25e-7}{\meter} \]

\textbf{Frequency:}
\[ \nu = \frac{c}{\lambda} = \frac{3.00\times10^{8}}{5.25\times10^{-7}}
   = \SI{5.71e14}{\per\second} \]

\textbf{Energy per photon:}
\[ E = h\nu = (6.626\times10^{-34})(5.71\times10^{14})
   = \SI{3.78e-19}{\joule} \]

\textbf{Per mole:}
\[ (3.78\times10^{-19})(6.022\times10^{23})
   = \SI{2.28e5}{\joule\per\mole} = \SI{228}{\kilo\joule\per\mole} \]

\textbf{Two sense checks worth thirty seconds.} A single photon carries
around $10^{-19}$ J --- if your answer is $10^{-30}$ or $10^{-5}$, a
conversion went wrong. And a mole of visible photons carries a couple of
hundred kJ, comparable to a chemical bond energy, which is precisely why
visible light can drive chemistry.
\end{workedexample}

\begin{yourturn}[YOUR TURN 1 --- four questions]
\begin{enumerate}[leftmargin=*,itemsep=0.5em]
  \item Frequency of light with $\lambda = \SI{650}{\nano\meter}$:
        \workspace[1.6cm]
  \item Energy of one photon of that light: \workspace[1.6cm]
  \item Wavelength, in nm, of a photon carrying
        \SI{4.50e-19}{\joule}: \workspace[1.8cm]
  \item Which has more energy per photon, a \SI{200}{\nano\meter}
        ultraviolet photon or a \SI{700}{\nano\meter} red one? Answer
        without calculating. \workspace[1.4cm]
\end{enumerate}
\selfcheck{(a) \SI{4.62e14}{\per\second} \quad (b)
\SI{3.06e-19}{\joule} \quad (c) \SI{442}{\nano\meter} \quad (d) the
ultraviolet one --- shorter wavelength}
\keyonly{\begin{apcanswer}(a) $\nu = 3.00\times10^{8}/6.50\times10^{-7}
= \mathbf{4.62\times10^{14}}$~s$^{-1}$.
(b) $E = (6.626\times10^{-34})(4.62\times10^{14}) =
\mathbf{3.06\times10^{-19}}$~J.
(c) $\lambda = hc/E = (6.626\times10^{-34})(3.00\times10^{8}) /
4.50\times10^{-19} = 4.42\times10^{-7}$~m $= \mathbf{442}$~nm.
(d) The \textbf{ultraviolet} photon. $E = hc/\lambda$, so energy is
inversely proportional to wavelength --- 200 nm is shorter, therefore
more energetic, which is why UV damages skin and red light does
not.\end{apcanswer}}
\end{yourturn}

% =====================================================================
\section*{\sffamily Ladder 2 \textbullet\ The photoelectric effect}
\zsec{7.2}

Shine light on a metal and electrons are ejected --- but only if the
light's \textbf{frequency} exceeds a threshold. Brighter light below that
threshold ejects nothing, however long you wait.

That single observation is what forced the particle picture. A wave would
let energy accumulate, so any colour should eventually work. A stream of
photons would not: one photon hits one electron, and if that photon is
too weak, nothing happens no matter how many arrive.

\[ E_{\text{photon}} = \Phi + \mathrm{KE}_{\text{electron}} \]
where $\Phi$ (the \term{work function}) is the minimum energy to free an
electron. Energy above the work function becomes kinetic energy.

\begin{workedexample}[I do: reading the energy budget]
A metal has a work function of \SI{3.20e-19}{\joule}. Light of
wavelength \SI{400.}{\nano\meter} strikes it.

\textbf{Photon energy:}
\[ E = \frac{hc}{\lambda}
   = \frac{(6.626\times10^{-34})(3.00\times10^{8})}{4.00\times10^{-7}}
   = \SI{4.97e-19}{\joule} \]

\textbf{Is it above the threshold?} $4.97 > 3.20$, so yes --- electrons
are ejected.

\textbf{Kinetic energy of each ejected electron:}
\[ \mathrm{KE} = 4.97\times10^{-19} - 3.20\times10^{-19}
   = \SI{1.77e-19}{\joule} \]

\textbf{Now the conceptual half.} What if we doubled the
\emph{intensity} at the same wavelength? Twice as many photons arrive, so
twice as many electrons are ejected --- but each still receives exactly
one photon's worth of energy, so the kinetic energy per electron is
\textbf{unchanged}. Intensity controls how many; frequency controls how
energetic. Confusing those two is the classic error here.
\end{workedexample}

\begin{yourturn}[YOUR TURN 2 --- four questions]
A metal has work function \SI{4.00e-19}{\joule}.
\begin{enumerate}[leftmargin=*,itemsep=0.5em]
  \item Energy of a \SI{450}{\nano\meter} photon: \workspace[1.6cm]
  \item Are electrons ejected, and if so with what kinetic energy?
        \workspace[1.6cm]
  \item What is the longest wavelength that will eject any electron from
        this metal? \workspace[1.8cm]
  \item Very intense \SI{600}{\nano\meter} light is shone on the metal
        for an hour. Predict what happens, and why. \workspace[1.6cm]
\end{enumerate}
\selfcheck{(a) \SI{4.42e-19}{\joule} \quad (b) yes, KE $=$
\SI{4.2e-20}{\joule} \quad (c) \SI{497}{\nano\meter} \quad (d) nothing
--- each photon is below the work function}
\keyonly{\begin{apcanswer}(a) $E = (6.626\times10^{-34})
(3.00\times10^{8})/4.50\times10^{-7} =
\mathbf{4.42\times10^{-19}}$~J. (b) $4.42 > 4.00$, so \textbf{yes};
KE $= (4.42 - 4.00)\times10^{-19} = \mathbf{4.2\times10^{-20}}$~J.
(c) At threshold all the photon energy goes to freeing the electron:
$\lambda = hc/\Phi = (6.626\times10^{-34})(3.00\times10^{8})/
4.00\times10^{-19} = 4.97\times10^{-7}$~m $= \mathbf{497}$~nm.
(d) \textbf{Nothing at all.} A 600 nm photon carries
$3.31\times10^{-19}$~J, below the $4.00\times10^{-19}$~J work function.
Intensity supplies more photons, not stronger ones, and one electron
absorbs one photon --- so no electron ever gathers enough energy,
however long the light shines.\end{apcanswer}}
\end{yourturn}

% =====================================================================
\section*{\sffamily Ladder 3 \textbullet\ Line spectra and the Bohr model}
\zsec{7.3--7.4}

Heat hydrogen and it emits only \emph{certain} wavelengths --- a line
spectrum, not a continuous rainbow. That is the direct evidence that
electron energies are \textbf{quantized}: an electron can occupy only
allowed levels, and a photon is emitted with exactly the energy of the
gap it falls across.

For hydrogen only:
\[ E_n = -2.178\times10^{-18}\left(\frac{1}{n^{2}}\right)~\text{J}
   \qquad
   \Delta E = E_{\text{final}} - E_{\text{initial}} \]

The negative sign means bound. $n = \infty$ is zero energy --- the
electron just barely free --- so every bound level lies below it.

\begin{workedexample}[I do: an emission line, end to end]
Find the wavelength emitted when a hydrogen electron falls from
$n = 4$ to $n = 2$.

\textbf{The two levels:}
\[ E_4 = \frac{-2.178\times10^{-18}}{16} = \SI{-1.361e-19}{\joule}
   \qquad
   E_2 = \frac{-2.178\times10^{-18}}{4} = \SI{-5.445e-19}{\joule} \]

\textbf{The change:}
\[ \Delta E = E_2 - E_4 = -5.445\times10^{-19} - (-1.361\times10^{-19})
   = \SI{-4.084e-19}{\joule} \]

Negative because the atom \emph{lost} energy. The photon carries that
energy away, so use its magnitude:

\[ \lambda = \frac{hc}{|\Delta E|}
   = \frac{(6.626\times10^{-34})(3.00\times10^{8})}{4.084\times10^{-19}}
   = 4.87\times10^{-7}~\text{m} = \mathbf{487~nm} \]

\textbf{Check against reality:} 487 nm is blue-green, and the
$4 \to 2$ line is indeed one of the visible hydrogen lines. Any
transition ending at $n = 2$ lands in or near the visible; those ending
at $n = 1$ are far more energetic and fall in the ultraviolet.
\end{workedexample}

\begin{yourturn}[YOUR TURN 3 --- four questions]
\begin{enumerate}[leftmargin=*,itemsep=0.5em]
  \item Energy of the $n = 3$ level in hydrogen: \workspace[1.4cm]
  \item $\Delta E$ for an electron falling from $n = 3$ to $n = 2$:
        \workspace[1.6cm]
  \item Wavelength of the photon emitted in (b): \workspace[1.8cm]
  \item Is energy absorbed or emitted for $n = 1 \to n = 3$, and how do
        you know from the sign? \workspace[1.6cm]
\end{enumerate}
\selfcheck{(a) \SI{-2.420e-19}{\joule} \quad (b)
\SI{-3.025e-19}{\joule} \quad (c) \SI{657}{\nano\meter} \quad
(d) absorbed --- $\Delta E$ is positive}
\keyonly{\begin{apcanswer}(a) $E_3 = -2.178\times10^{-18}/9 =
\mathbf{-2.420\times10^{-19}}$~J. (b) $\Delta E = E_2 - E_3 =
-5.445\times10^{-19} + 2.420\times10^{-19} =
\mathbf{-3.025\times10^{-19}}$~J. (c) $\lambda = hc/|\Delta E| =
1.988\times10^{-25}/3.025\times10^{-19} = 6.57\times10^{-7}$~m
$= \mathbf{657}$~nm --- the red hydrogen line. (d) \textbf{Absorbed.}
Going up in $n$ makes $\Delta E$ \textbf{positive}: the atom gains
energy, which it must take from a photon. Emission always gives a
negative $\Delta E$.\end{apcanswer}}
\end{yourturn}

% =====================================================================
\section*{\sffamily Ladder 4 \textbullet\ Quantum numbers}
\zsec{7.6--7.8}

Four numbers address one electron, and each answers a different question.

\begin{center}
\renewcommand{\arraystretch}{1.25}
\begin{tabular}{@{}p{1.4cm}p{3.4cm}p{4.2cm}p{4.2cm}@{}}
\hline
\textbf{Symbol} & \textbf{Name} & \textbf{Allowed values} &
  \textbf{What it fixes} \\
\hline
$n$ & principal & $1, 2, 3, \ldots$ & shell; size and energy \\
$\ell$ & angular momentum & $0$ to $n-1$ & subshell shape
  ($0=s$, $1=p$, $2=d$, $3=f$) \\
$m_\ell$ & magnetic & $-\ell$ to $+\ell$ & which orbital, i.e.\
  orientation \\
$m_s$ & spin & $+\tfrac{1}{2}$ or $-\tfrac{1}{2}$ & which of the two
  electrons \\
\hline
\end{tabular}
\end{center}

The \term{Pauli exclusion principle} says no two electrons in an atom
share all four --- which is exactly why an orbital holds two electrons
and no more.

\begin{workedexample}[I do: counting, then judging legality]
\textbf{(a) How many orbitals are in the $n = 3$ shell, and how many
electrons can it hold?}

$\ell$ runs $0, 1, 2$ --- so $3s$, $3p$, $3d$. Counting orbitals by
$m_\ell$: $s$ has 1, $p$ has 3, $d$ has 5, total $\mathbf{9}$ orbitals.
Two electrons each gives $\mathbf{18}$ --- and $2n^{2} = 2(9) = 18$
confirms it.

\textbf{(b) Is $n = 2$, $\ell = 2$ legal?} No. $\ell$ may run only to
$n - 1 = 1$, so there is no $2d$ subshell --- the first $d$ subshell is
$3d$.

\textbf{(c) Is $n = 3$, $\ell = 1$, $m_\ell = -2$ legal?} No. With
$\ell = 1$, $m_\ell$ is limited to $-1, 0, +1$; the value $-2$ lies
outside that range.

\textbf{The procedure:} check the numbers \emph{in order}. $\ell$ is
constrained by $n$, and $m_\ell$ is constrained by $\ell$ --- so an
illegal set usually breaks at the first constraint you test.
\end{workedexample}

\begin{yourturn}[YOUR TURN 4 --- four questions]
\begin{enumerate}[leftmargin=*,itemsep=0.5em]
  \item How many orbitals and how many electrons in the $n = 4$ shell?
        \workspace[1.6cm]
  \item Which subshell is $n = 4$, $\ell = 2$, and how many orbitals does
        it contain? \workspace[1.4cm]
  \item Is $n = 2$, $\ell = 1$, $m_\ell = 0$,
        $m_s = +\tfrac{1}{2}$ a legal set? \workspace[1.4cm]
  \item Explain, using the Pauli principle, why an orbital holds at most
        two electrons. \workspace[1.6cm]
\end{enumerate}
\selfcheck{(a) 16 orbitals, 32 electrons \quad (b) $4d$, 5 orbitals
\quad (c) yes \quad (d) only two $m_s$ values exist, so a third electron
would duplicate all four numbers}
\keyonly{\begin{apcanswer}(a) $\ell = 0,1,2,3$ gives $1+3+5+7 =
\mathbf{16}$ orbitals and $\mathbf{32}$ electrons ($2n^{2} = 32$).
(b) $\ell = 2$ is $d$, so \textbf{$4d$}, with $m_\ell = -2,-1,0,1,2$ ---
\textbf{5} orbitals. (c) \textbf{Yes} --- $\ell = 1 \le n-1$,
$m_\ell = 0$ lies within $-1$ to $+1$, and $m_s$ is allowed. This is a
$2p$ electron. (d) An orbital is one fixed set of $n$, $\ell$,
$m_\ell$. The only remaining freedom is $m_s$, which has just two
values, so two electrons can differ. A third would have to repeat all
four numbers, which Pauli forbids.\end{apcanswer}}
\end{yourturn}

% =====================================================================
\section*{\sffamily Ladder 5 \textbullet\ Electron configurations}
\zsec{7.11}

Fill in order of increasing energy, obeying three rules:
\begin{itemize}[leftmargin=*,itemsep=0.15em]
  \item \term{Aufbau}: lowest-energy subshell first.
  \item \term{Pauli}: two electrons per orbital, opposite spins.
  \item \term{Hund}: within a subshell, occupy every orbital singly
        before pairing --- electrons repel, so they spread out first.
\end{itemize}

The periodic table \emph{is} the filling order: read across the rows and
the blocks give you the sequence without memorizing a diagonal diagram.

\begin{workedexample}[I do: an atom, an ion, and an anomaly]
\textbf{(a) Sulfur (Z = 16).}
\[ 1s^{2}\,2s^{2}\,2p^{6}\,3s^{2}\,3p^{4}
   \qquad\text{or}\qquad [\ce{Ne}]\,3s^{2}\,3p^{4} \]
Check the electron count: $2+2+6+2+4 = 16$. Always run that sum.

\textbf{(b) The \ce{S^2-} ion.} Two extra electrons complete the $3p$:
\[ [\ce{Ne}]\,3s^{2}\,3p^{6} = [\ce{Ar}] \]

\textbf{(c) Copper (Z = 29), an anomaly.} Aufbau predicts
$[\ce{Ar}]\,4s^{2}\,3d^{9}$, but the actual configuration is
\[ [\ce{Ar}]\,4s^{1}\,3d^{10} \]
A filled $d$ subshell is low enough in energy to be worth promoting one
$4s$ electron. Chromium does the same thing for a half-filled shell:
$[\ce{Ar}]\,4s^{1}\,3d^{5}$. These two are the anomalies worth knowing.

\textbf{The trap in (b) that catches everyone:} for a transition-metal
\emph{cation}, electrons leave the $4s$ \emph{before} the $3d$, even
though $4s$ filled first. \ce{Fe^2+} is $[\ce{Ar}]\,3d^{6}$, not
$[\ce{Ar}]\,4s^{2}\,3d^{4}$.
\end{workedexample}

\begin{yourturn}[YOUR TURN 5 --- four questions]
\begin{enumerate}[leftmargin=*,itemsep=0.5em]
  \item Write the full configuration for calcium (Z = 20).
        \workspace[1.6cm]
  \item Write the noble-gas configuration for bromine (Z = 35).
        \workspace[1.6cm]
  \item Write the configuration for \ce{Fe^3+} (Fe is Z = 26).
        \workspace[1.6cm]
  \item How many unpaired electrons does nitrogen (Z = 7) have, and
        which rule decides it? \workspace[1.6cm]
\end{enumerate}
\selfcheck{(a) $1s^{2}2s^{2}2p^{6}3s^{2}3p^{6}4s^{2}$ \quad (b)
$[\ce{Ar}]4s^{2}3d^{10}4p^{5}$ \quad (c) $[\ce{Ar}]3d^{5}$ \quad
(d) three, by Hund's rule}
\keyonly{\begin{apcanswer}(a) $\mathbf{1s^{2}2s^{2}2p^{6}3s^{2}3p^{6}
4s^{2}}$ --- the counts sum to 20. (b) $\mathbf{[\ce{Ar}]4s^{2}3d^{10}
4p^{5}}$: 18 from argon plus $2+10+5 = 17$ gives 35.
(c) Iron is $[\ce{Ar}]4s^{2}3d^{6}$. Remove the two $4s$ electrons
\emph{first}, then one $3d$: $\mathbf{[\ce{Ar}]3d^{5}}$ --- a
half-filled $d$ subshell, which is part of why \ce{Fe^3+} is so common.
(d) Nitrogen is $1s^{2}2s^{2}2p^{3}$. By \textbf{Hund's rule} the three
$2p$ electrons occupy the three $2p$ orbitals singly with parallel
spins, giving \textbf{three} unpaired electrons.\end{apcanswer}}
\end{yourturn}

% =====================================================================
\section*{\sffamily Ladder 6 \textbullet\ Photoelectron spectroscopy}
\zsec{7.9}

A PES spectrum is a direct picture of an atom's subshells. Read it with
two rules and nothing else:

\begin{center}
\fbox{\parbox{0.92\linewidth}{\centering
\textbf{Peak position} $=$ binding energy $=$ how tightly that subshell
is held.\\[2pt]
\textbf{Peak area (height)} $=$ \emph{how many electrons} are in that
subshell.}}
\end{center}

\begin{apctrap}
Do not confuse this with a \textbf{mass spectrum}, where height gives
isotope \emph{abundance}. In PES the vertical axis counts
\emph{electrons in a subshell}, and the ratios come out as
$2:2:6:2:\ldots$ --- exactly the subshell populations of a
configuration.
\end{apctrap}

\begin{workedexample}[I do: identifying an element from its spectrum]
A PES spectrum shows peaks at 1.36, 0.10, and 0.06 MJ/mol with relative
areas $2 : 2 : 1$.

\textbf{Read the areas as electron counts.} $2 : 2 : 1$ means 2, 2, and
1 electrons --- five in total, so $Z = 5$: \textbf{boron}.

\textbf{Assign the peaks.} The largest binding energy belongs to the
electrons closest to the nucleus and least shielded:
\[ 1.36 = 1s^{2} \qquad 0.10 = 2s^{2} \qquad 0.06 = 2p^{1} \]
Configuration $1s^{2}2s^{2}2p^{1}$ --- boron, confirmed twice over.

\textbf{Why 1.36 is so much larger than the rest.} The $1s$ electrons
sit closest to the $+5$ nucleus with essentially nothing between them
and it. The $n = 2$ electrons are further out \emph{and} shielded by
the $1s$ pair, so the nucleus pulls on them far more weakly --- a
Coulombic argument, and the only kind that earns credit.

\textbf{And why $2s$ exceeds $2p$:} the $2s$ orbital penetrates closer
to the nucleus, so a $2s$ electron spends more time inside the shielding
and is held more tightly.
\end{workedexample}

\begin{yourturn}[YOUR TURN 6 --- four questions]
A PES spectrum shows peaks at 6.84, 0.68, 0.40 MJ/mol with relative
areas $2 : 2 : 4$.
\begin{enumerate}[leftmargin=*,itemsep=0.5em]
  \item How many electrons in total, and which element is it?
        \workspace[1.4cm]
  \item Assign each peak to a subshell. \workspace[1.6cm]
  \item Explain, in Coulombic terms, why the 6.84 peak is so much larger
        than the others. \workspace[1.8cm]
  \item In a PES spectrum, what does peak \emph{height} represent --- and
        what does it represent in a mass spectrum? \workspace[1.6cm]
\end{enumerate}
\selfcheck{(a) 8 electrons --- oxygen \quad (b) $1s^{2}$, $2s^{2}$,
$2p^{4}$ \quad (c) closest to the nucleus and unshielded \quad
(d) electrons per subshell; isotope abundance}
\keyonly{\begin{apcanswer}(a) $2+2+4 = \mathbf{8}$ electrons, so
$Z = 8$: \textbf{oxygen}. (b) $6.84 = \mathbf{1s^{2}}$,
$0.68 = \mathbf{2s^{2}}$, $0.40 = \mathbf{2p^{4}}$. (c) The $1s$
electrons are \textbf{closest to the nucleus} and have \textbf{no
shielding} between them and its $+8$ charge, so the Coulombic attraction
is enormous. The $n = 2$ electrons are both further away and shielded by
the $1s$ pair, so far less energy frees them. (d) In \textbf{PES}, height
(area) gives the \textbf{number of electrons in that subshell}. In a
\textbf{mass spectrum} it gives the \textbf{relative abundance of an
isotope}. Same-looking axis, entirely different
quantity.\end{apcanswer}}
\end{yourturn}

% =====================================================================
\section*{\sffamily Ladder 7 \textbullet\ Successive ionization energies}
\zsec{7.12}

Removing electrons one at a time gives IE$_1$, IE$_2$, IE$_3$, \ldots ---
each larger than the last, because pulling a negative electron away from
an increasingly positive ion is harder every time.

The information is in the \textbf{jump}. A sudden large increase means
the next electron had to come from a \emph{new, inner} shell --- much
closer to the nucleus and much less shielded. Count the electrons removed
before the jump and you have the number of valence electrons, hence the
group.

\begin{workedexample}[I do: identifying a group from the jumps]
An element has successive ionization energies, in MJ/mol, of
\[ 0.74,\quad 1.45,\quad 7.73,\quad 10.5 \]
Which group does it belong to?

\textbf{Look at the ratios, not the differences.}
\[ \frac{1.45}{0.74} = 2.0 \qquad
   \frac{7.73}{1.45} = 5.3 \qquad
   \frac{10.5}{7.73} = 1.4 \]

\textbf{The jump is from IE$_2$ to IE$_3$} --- a factor of more than
five, against modest steps either side.

\textbf{Interpretation.} Two electrons came off comparatively easily;
the third was drastically harder. So the atom has \textbf{two valence
electrons} and is in \textbf{group 2A}. (The values are magnesium's.)

\textbf{Say why in Coulombic terms.} The first two electrons came from
the $3s$ subshell --- far from the nucleus and shielded by the $n = 1$
and $n = 2$ electrons. The third had to come from $n = 2$: much closer
to the nucleus and shielded only by the $1s$ pair, so the attraction to
overcome is several times greater.

\textbf{What earns zero here:} ``the ion had a full octet and wanted to
stay stable.'' That names the pattern instead of explaining it. Every
ionization-energy explanation must cite nuclear charge, distance, or
shielding.
\end{workedexample}

\begin{yourturn}[YOUR TURN 7 --- four questions]
An element has IEs, in MJ/mol, of $1.01$, $1.91$, $2.91$, $4.96$,
$6.27$, $21.3$, $25.4$.
\begin{enumerate}[leftmargin=*,itemsep=0.5em]
  \item Between which two ionizations is the large jump?
        \workspace[1.4cm]
  \item How many valence electrons, and which group? \workspace[1.4cm]
  \item Explain the jump in Coulombic terms --- no octet language.
        \workspace[2cm]
  \item Why is IE$_2$ always larger than IE$_1$ for any element?
        \workspace[1.6cm]
\end{enumerate}
\selfcheck{(a) between IE$_5$ and IE$_6$ \quad (b) five valence
electrons, group 5A \quad (c) the sixth comes from an inner shell ---
closer, less shielded \quad (d) removing from a more positive ion}
\keyonly{\begin{apcanswer}(a) Between \textbf{IE$_5$ and IE$_6$}:
$21.3/6.27 = 3.4$, against steps of at most 1.9 elsewhere.
(b) \textbf{Five} valence electrons, so \textbf{group 5A} --- these are
phosphorus's values. (c) The first five electrons came from the $n = 3$
shell, relatively far from the nucleus and shielded by the ten $n = 1$
and $n = 2$ electrons. The sixth had to come from $n = 2$: much
\textbf{closer to the nucleus} and shielded only by the $1s$ pair, so the
Coulombic attraction to overcome is several times larger. (d) Because
IE$_2$ removes an electron from an ion that already carries a $1+$
charge. The same nuclear charge now holds fewer electrons, so each
remaining one is attracted more strongly and is harder to
remove.\end{apcanswer}}
\end{yourturn}

% =====================================================================
\section*{\sffamily Ladder 8 \textbullet\ Periodic trends, one engine}
\zsec{7.12--7.13}

Three trends, and they all run on the same Coulombic reasoning:

\begin{center}
\renewcommand{\arraystretch}{1.25}
\begin{tabular}{@{}p{3.2cm}p{4.4cm}p{5.6cm}@{}}
\hline
\textbf{Trend} & \textbf{Across a period ($\to$)} &
  \textbf{Down a group ($\downarrow$)} \\
\hline
Atomic radius & decreases & increases \\
Ionization energy & increases & decreases \\
Electronegativity & increases & decreases \\
\hline
\end{tabular}
\end{center}

\textbf{Why, in one paragraph.} Across a period, protons are added while
electrons enter the \emph{same} shell, so shielding barely changes: the
effective nuclear charge rises, the electrons are pulled in tighter, and
the atom shrinks while holding its electrons more strongly. Down a
group, each element adds a whole new shell: the outer electrons are much
\emph{further} from the nucleus and shielded by every inner shell, so
they are held far more loosely.

\begin{apctrap}
\textbf{The language rule for this entire chapter.} Every trend
explanation must cite \textbf{nuclear charge, distance, or shielding}.
Answers built on ``full shells'', ``octets'', or ``wanting to be
stable'' describe the pattern rather than explaining it, and earn
\textbf{zero} on this course's rubrics --- however confidently they are
written.
\end{apctrap}

\begin{workedexample}[I do: comparisons, each justified Coulombically]
\textbf{(a) Which has the larger radius, \ce{Na} or \ce{Mg}?}
\textbf{\ce{Na}}. Magnesium has one more proton pulling on electrons in
the same $n = 3$ shell with essentially the same shielding, so its
effective nuclear charge is higher and its electron cloud is drawn in
tighter.

\textbf{(b) Which has the higher first ionization energy, \ce{Li} or
\ce{Cs}?} \textbf{\ce{Li}}. Its outer electron sits in $n = 2$, close to
the nucleus and shielded only by the $1s$ pair. Caesium's outer electron
is in $n = 6$, far away and shielded by five inner shells, so it is held
weakly and easily removed --- which is exactly why caesium is so
reactive.

\textbf{(c) Which is larger, \ce{Cl} or \ce{Cl-}?} \textbf{\ce{Cl-}}.
The added electron does not change the nuclear charge, but it increases
electron--electron repulsion in the valence shell, so the cloud expands.
Anions are always larger than their parent atoms; cations are always
smaller, for the mirror-image reason.

\textbf{Notice what none of these answers mentioned:} octets, full
shells, or stability. Each one named a charge, a distance, or a
shielding difference.
\end{workedexample}

\begin{yourturn}[YOUR TURN 8 --- four questions]
Answer each with a Coulombic justification.
\begin{enumerate}[leftmargin=*,itemsep=0.5em]
  \item Larger radius: \ce{K} or \ce{Br}? \workspace[1.6cm]
  \item Higher first ionization energy: \ce{Ne} or \ce{Na}?
        \workspace[1.6cm]
  \item Larger: \ce{Mg} or \ce{Mg^2+}? \workspace[1.6cm]
  \item \ce{Na+} and \ce{F-} both have 10 electrons. Which is smaller,
        and why? \workspace[1.8cm]
\end{enumerate}
\selfcheck{(a) \ce{K} \quad (b) \ce{Ne} \quad (c) \ce{Mg} \quad
(d) \ce{Na+} --- more protons pulling on the same ten electrons}
\keyonly{\begin{apcanswer}(a) \textbf{\ce{K}}. Both are in period 4, but
bromine has far more protons acting on electrons in the same shell with
similar shielding, so its higher effective nuclear charge contracts it.
(b) \textbf{\ce{Ne}}. Its outer electron is in $n = 2$, close in and
shielded only by $1s$; sodium's is in $n = 3$, further out and shielded
by all ten inner electrons, so it is held much more weakly.
(c) \textbf{\ce{Mg}}. Removing two electrons empties the $3s$ subshell
entirely, so the ion's outermost electrons are now $n = 2$ --- a whole
shell closer --- and the unchanged $+12$ nuclear charge now pulls on
only ten electrons. (d) \textbf{\ce{Na+}}. The two are isoelectronic, so
shielding is identical, but sodium has \textbf{11 protons} against
fluorine's \textbf{9}. More nuclear charge on the same ten electrons
means a stronger pull and a smaller radius.\end{apcanswer}}
\end{yourturn}

% =====================================================================
\section*{\sffamily Mastery tracker}

Tick a row only if \textbf{all four} YOUR TURN questions were right on
the first attempt.

\renewcommand{\arraystretch}{1.35}
\noindent\begin{tabular}{@{}c p{6.4cm} c p{4.6cm}@{}}
\toprule
\textbf{First try?} & \textbf{Skill} & \textbf{Ladder} &
  \textbf{If not, re-read\ldots}\\
\midrule
$\square$ & wavelength, frequency, energy & 1 & convert to metres first \\
$\square$ & the photoelectric effect & 2 & intensity vs.\ frequency \\
$\square$ & line spectra and Bohr levels & 3 & sign of $\Delta E$ \\
$\square$ & quantum numbers & 4 & check them in order \\
$\square$ & electron configurations & 5 & $4s$ leaves before $3d$ \\
$\square$ & photoelectron spectroscopy & 6 & area $=$ electron count \\
$\square$ & successive ionization energies & 7 & find the jump \\
$\square$ & periodic trends & 8 & charge, distance, shielding \\
\bottomrule
\end{tabular}

\begin{apcnote}[Scoring yourself honestly]
8/8: this chapter feeds Units 1 and 2 directly, so you are in good shape
for both. 6--7: redo the missed ladders tomorrow, not today.
5 or fewer: check where the misses fall. Ladders 1--3 are
\emph{arithmetic} --- almost always a unit conversion, usually nm to m.
Ladders 6--8 are \emph{explanation}, and there the failure is nearly
always reaching for octets and stability instead of naming a nuclear
charge, a distance, or a shielding difference. Those are two different
problems with two different fixes; diagnose which one you have before
redoing anything.
\end{apcnote}

\end{document}
